\documentclass[12pt,a4paper]{report}
\usepackage[margin=25mm]{geometry}
\usepackage{fontspec}\setmainfont{TeX Gyre Termes}
\usepackage{graphicx,amsmath,longtable,booktabs,array,hyperref,xurl}
\hypersetup{hidelinks}
\setcounter{tocdepth}{1}
\begin{document}

\begin{titlepage}\centering K. K. Wagh Institute of Engineering Education \& Research\par\vspace{2cm}{\Large Explainable AI Framework with Explanation Validation for Financial Decision Support}\par\vspace{1cm}TrustFin AI\par Group Project Report\par\vspace{1cm}Ashwin Patil (46), Aaditya Patil (44), Kanad Nerkar (41), Ved Kalantri (28)\par Project Guide: Prof. Rupali M. Bora\par Academic Year 2026-2027\end{titlepage}
\chapter*{Certificate}\addcontentsline{toc}{chapter}{Certificate}
This group project report entitled Explainable AI Framework with Explanation Validation for Financial Decision Support, developed as TrustFin AI, is submitted by Ashwin Patil (46), Aaditya Patil (44), Kanad Nerkar (41) and Ved Kalantri (28), Department of Information Technology, in partial fulfilment of the requirements for the degree of Bachelor of Technology.

The report records work carried out under the guidance of Prof. Rupali M. Bora during academic year 2026-2027. Institutional certification, approval and signatures are reserved for the institute; this draft does not represent completed certification.

Project Guide: Prof. Rupali M. Bora

Head of Department: \_\_\_\_\_\_\_\_\_\_\_\_\_\_\_\_\_\_\_\_

Director: \_\_\_\_\_\_\_\_\_\_\_\_\_\_\_\_\_\_\_\_

Date: \_\_\_\_\_\_\_\_\_\_\_\_\_\_\_\_\_\_\_\_    Place: Nashik

\chapter*{Acknowledgement}\addcontentsline{toc}{chapter}{Acknowledgement}
The group expresses its sincere gratitude to Prof. Rupali M. Bora for guidance and feedback during the development of TrustFin AI. We thank the Department of Information Technology and K. K. Wagh Institute of Engineering Education \& Research for the academic environment in which the project was carried out.

The project brought together data preparation, machine learning, explainability, backend services, database persistence and interface development. We acknowledge the authors and research communities whose published work on credit scoring, SHAP, LIME, explanation evaluation and responsible machine learning informed the study. Relevant publications are cited in the bibliography.

We also thank the project group for the shared effort involved in designing, implementing, reviewing and documenting the system.

\chapter*{Abstract}\addcontentsline{toc}{chapter}{Abstract}
TrustFin AI implements an explainable loan decision-support prototype that combines product-specific classification with local explanations and quantitative diagnostic checks. The framework addresses a practical difficulty in automated financial assessment: a prediction may be difficult to interpret, while a plausible explanation does not by itself establish fidelity, robustness or fairness.

The prototype supports home-loan and personal-loan applications through separate XGBoost models and preprocessing pipelines. Engineered features include amortized monthly installment and fixed-obligation-to-income ratio, with loan-to-value ratio for home loans. TreeSHAP and tabular LIME produce local feature contributions. The validation stage measures SHAP masking sensitivity, seeded SHAP stability, top-three cross-explainer overlap and individual gender-flip probability sensitivity. Their weighted combination forms a project-defined Trust Score, kept separate from predicted-class probability and affordability assessment. React/FastAPI interfaces, SQLAlchemy persistence, batch processing, history, recommendations and ReportLab reports complete the workflow.

Each product dataset contains 3,000 synthetic records generated by repository scripts. Saved home-loan metrics are 97.00\% accuracy and 0.9978 ROC-AUC; personal-loan metrics are 97.33\% and 0.9963. Accuracy and ROC-AUC were reproduced using the saved artifacts and deterministic test split. These results concern generated approval labels rather than observed repayment outcomes. Selected backend tests passed, while frontend and additional batch checks included failures retained in the evidence. The prototype reports solvency alongside the classifier result rather than enforcing a class override.

The contribution is an integrated, inspectable prediction and explanation-review workflow. The Trust Score is an engineering diagnostic, not an independently validated guarantee of lending quality. The report presents original conceptual design diagrams, concrete prototype implementation and measured verification as distinct but connected views.

Keywords: Explainable AI; XGBoost; SHAP; LIME; explanation validation; loan decision support; Trust Score.

\chapter*{Contents}\addcontentsline{toc}{chapter}{Contents}
\tableofcontents
\chapter*{List of Abbreviations}\addcontentsline{toc}{chapter}{List of Abbreviations}
\begin{longtable}{p{0.46\linewidth}p{0.46\linewidth}}

Abbreviation & Meaning\\
AI & Artificial Intelligence\\
API & Application Programming Interface\\
AUC & Area Under the ROC Curve\\
CIBIL & Credit Information Bureau (India) Limited; used here as a credit-score field\\
CSV & Comma-Separated Values\\
DFD & Data Flow Diagram\\
DTI & Debt-to-Income Ratio\\
EMI & Equated Monthly Instalment\\
FOIR & Fixed Obligation to Income Ratio\\
HTTP & Hypertext Transfer Protocol\\
JSON & JavaScript Object Notation\\
LIME & Local Interpretable Model-agnostic Explanations\\
LTV & Loan-to-Value Ratio\\
ML & Machine Learning\\
ORM & Object-Relational Mapper\\
PDF & Portable Document Format\\
ROC & Receiver Operating Characteristic\\
SHAP & SHapley Additive exPlanations\\
SQL & Structured Query Language\\
XAI & Explainable Artificial Intelligence\\
XGBoost & eXtreme Gradient Boosting\\
\end{longtable}
\chapter*{List of Figures}\addcontentsline{toc}{chapter}{List of Figures}
\listoffigures
\chapter*{List of Graphs}\addcontentsline{toc}{chapter}{List of Graphs}
\noindent Graphs are listed with figures using the same figure numbering.
\chapter*{List of Tables}\addcontentsline{toc}{chapter}{List of Tables}
\listoftables
\chapter{Introduction}
\section{Introduction}
TrustFin AI is a group project that studies how local explanations and explanation checks can be included in a loan decision-support application. The system accepts an applicant profile, selects a home-loan or personal-loan model, returns a binary approval prediction and probability, and generates SHAP and LIME explanations. Four validation values are then calculated and combined into a composite Trust Score. Results are presented through a browser dashboard and can be stored for later review.

The central concern is not simply whether a classifier can produce a high score on a test split. In a consequential setting such as credit assessment, users also need to understand what information influenced an output and whether a generated explanation behaves consistently under defined checks. TrustFin AI therefore treats explanation quality as a measurable part of the workflow. It does not treat the resulting Trust Score as proof that a lending decision is correct, fair, or suitable for deployment.

\section{Background}
Credit scoring uses applicant and account information to estimate the likelihood of a credit outcome. Statistical classification has been used for consumer credit for decades, with studies comparing a range of classification methods and discussing domain-specific problems such as incomplete applicant histories and reject inference [1, 2]. Modern ensemble methods can model non-linear interactions in structured records, but their internal logic is less direct to communicate than a small scorecard or decision tree.

An explanation method attempts to summarize a model’s behavior for a particular input. SHAP assigns feature contributions using a Shapley-value framework, while LIME fits an interpretable approximation around an instance. These methods answer related but not identical questions and can disagree. Moreover, an explanation that looks plausible is not necessarily faithful to the model, robust to small input changes, or neutral with respect to a protected feature. Research has therefore examined explanation quality as a separate evaluation problem rather than assuming that an explanation is valid merely because it was generated [3, 4].

\section{AI in Financial Decision-Making}
Structured credit data are a natural setting for supervised classification: records contain numerical fields such as income, requested amount, repayment term, and credit indicators, together with categorical attributes. TrustFin AI uses XGBoost for its two product-specific models. The model provides a decision and a probability for the approval class. In the application, the probability associated with the selected class is also returned as a confidence representation.

Automated or model-assisted finance requires attention to the complete decision path. Input handling affects what the model receives; preprocessing affects feature scale and representation; a prediction probability is not the same as an explanation; and an explanation is not an audit of model fairness. A useful decision-support interface should preserve these distinctions and show enough information for an officer or applicant to understand the result. General risk-management guidance likewise emphasizes that system risks depend on context and should be considered throughout development and use [5].

The product target in this prototype is a generated approval label, not a measured subsequent default event. The financial domain motivates the design, while the experiments evaluate the implemented classification and explanation workflow on synthetic records.

\section{Motivation}
The project was motivated by three connected needs. First, applicant-facing systems should communicate the factors associated with a model output in accessible language. Second, a single explainer may not be a sufficient basis for review, so the application should compare SHAP and LIME outputs and report differences. Third, explanation quality should be assessed along more than one dimension: whether influential features affect the model output, whether rankings persist after a small perturbation, whether two methods select similar leading features, and whether a protected-attribute counterfactual changes the probability.

The motivation is deliberately limited. The repository contains generated datasets and experimental model artifacts, not production lending records. TrustFin AI demonstrates an engineering workflow for prediction and explanation review; it is not a validated credit policy or a substitute for institutional model governance and human review.

The proposed framework brings prediction, explanation, validation and human review into one architecture. The local prototype realizes the central computational path through two product pipelines and a browser interface. The design diagrams communicate that broader organization; Chapter 5 identifies the concrete modules that execute it.

\section{Aim of the Project}
The project prototype delivers a loan decision-support system that produces productspecific predictions, presents local model explanations, and calculates defined checks on the generated explanations before displaying a composite explanation-reliability indicator.

\section{Objectives}
The project objectives are to:

1. construct separate, reproducible feature pipelines for home-loan and personal-loan applications;

2. train and evaluate XGBoost classifiers using fixed stratified train-test partitions;

3. integrate local TreeSHAP and tabular LIME explanations in a common response format;

4. calculate SHAP faithfulness and stability, cross-explainer feature consistency, and individual gender sensitivity;

5. combine the diagnostic values into the implemented weighted Trust Score while keeping predicted-class probability separate;

6. provide financial affordability assessment, applicant and officer presentations, stored history, batch processing and PDF reporting;

7. document observed prototype behavior, verification outcomes and the boundary between conceptual capabilities and implemented functions.

\section{Scope}
The implemented scope includes home-loan and personal-loan application forms, their trained models and engineered financial features, single-record prediction, SHAP and LIME explanations, four explanation checks, composite Trust Score, SQLite/PostgreSQL persistence, batch CSV processing, evaluation history, and PDF report generation. The interface provides applicant and bank-officer views; these are presentation modes selected in the client, not authenticated accounts.

The project does not implement a real lender integration, a production underwriting policy, secure user authentication, population-level fairness monitoring, probability calibration, or an independently validated Trust Score threshold. The interactive what-if calculator is a deterministic affordability scenario, not a live call to the prediction model. Fraud detection and unrelated financial products are outside the scope.

\section{Implemented Approach}
The implemented workflow is: collect a product-specific application; check required fields and types, normalize request fields; engineer repayment features; transform fields using the product preprocessor; obtain an XGBoost decision and probability; compute SHAP and LIME local explanations; evaluate four explanation properties; calculate the weighted Trust Score; return a structured response; and present, persist, or export the result. This sequence is expanded in Chapters 4 and 5.

\section{Organization of the Report}
Chapter 2 reviews credit classification, SHAP, LIME, explanation evaluation and fairness. Chapter 3 defines the problem, requirements and implementation boundaries. Chapter 4 presents the original conceptual design diagrams and maps them to the prototype. Chapter 5 documents datasets, financial calculations, model training, explanation algorithms, software modules and reporting. Chapter 6 reports the actual verification evidence, including failures and untested paths. Chapter 7 discusses saved model metrics and worked application cases. Chapter 8 retains the Git-supported development chronology. Chapter 9 concludes with the project contribution and realistic future scope. Appendices record the tested environment, source map and a retained evaluation example.

\chapter{Literature Survey}
\section{Introduction}
This chapter reviews work relevant to the implemented system: consumer credit classification, gradient-boosted trees, post-hoc local explanation, objective explanation checks, and fairness measures. The focus is on concepts used to motivate or interpret TrustFin AI. The project does not claim to introduce new versions of SHAP, LIME, or the fairness criteria in the literature.

\section{Machine Learning in Loan Approval}
Hand and Henley surveyed statistical classification methods used in consumer credit scoring and described issues specific to credit data, including the limited public availability of lending datasets [1]. Later benchmarking work compared a large set of classifiers across multiple credit-scoring datasets, showing why model choice should be assessed empirically rather than inferred from a single algorithm’s reputation [2]. These studies provide context for using supervised learning while also highlighting the importance of dataset properties and evaluation design.

Gradient boosting builds an additive model by fitting successive learners to reduce a loss function. Friedman formalized gradient boosting as a function-approximation procedure [6]. XGBoost provides an optimized tree-boosting implementation with regularization and systems techniques designed for efficient learning on structured data [7]. TrustFin AI uses XGBoost because it is suitable for its tabular, mixed-type loan records and supports the selected TreeSHAP workflow. This engineering choice alone does not establish superiority over other classifiers; no current product-specific comparison experiment is stored in the repository.

\section{Explainable Artificial Intelligence}
Explainable AI includes methods intended to make model behavior more understandable to people. Explanations may describe global patterns, local feature attributions, examplebased behavior, or simpler models. TrustFin AI implements local post-hoc explanations for individual loan records. Such explanations should be interpreted as descriptions of model behavior under a chosen method and representation. They do not, by themselves, establish a causal relationship between a feature and repayment or prove that changing a feature would change a real-world outcome.

\section{SHAP}
Lundberg and Lee presented SHAP as a unified framework for additive feature attribution, connecting local explanation methods to Shapley-value principles [8]. In a local additive explanation, the model output is represented through a baseline and feature contributions. The sum of contributions is associated with the explained prediction in the chosen output space. The exact interpretation depends on the explainer, model, and output convention.

TrustFin AI initializes a SHAP TreeExplainer for each loaded XGBoost product model. For a request, the transformed instance is passed to the explainer and returned contributions are converted into a common schema containing feature name, value, contribution, and direction. The dashboard uses these contributions to show leading decision drivers. In this implementation, the explainer operates on preprocessed model features while the display path attempts to retain applicant-level values where names align. This mapping should be reviewed when interpreting categorical one-hot features and engineered columns.

Tree-specific explanation research extends the additive-attribution perspective to efficient explanation of tree ensembles [12]. The implementation uses TreeExplainer for XGBoost, whereas the general SHAP paper [8] supplies the foundational additive framework. In this report, the algorithm name does not imply causal explanation or a probability contribution.

\section{LIME}
Ribeiro, Singh, and Guestrin introduced LIME as a model-agnostic approach that explains a prediction by fitting a simple interpretable model in the neighborhood of the instance [9]. LIME’s explanation is local: it depends on the perturbation distribution, similarity weighting, feature representation, and surrogate fitting settings. It is not a global account of the classifier.

TrustFin AI initializes a tabular LIME explainer with a transformed background sample and the model’s transformed feature names. At prediction time the model’s probability function is passed to LIME. The implementation returns feature contributions in the common explanation schema. Applicant-specific deterministic seeds are used to reduce run-to-run variation in LIME’s sampling for identical inputs. Seed control improves repeatability under this implementation; it does not remove sampling approximation or establish correctness of a local surrogate.

\section{Explanation Validation}
Explanation quality is multi-dimensional. A feature ranking can appear stable while being disconnected from model behavior; an explanation can change under small perturbations; and separate methods can select different features. Yeh et al. study measures related to infidelity and sensitivity and show that explanation evaluation involves choices about perturbations and objective functions [3]. Alvarez-Melis and Jaakkola identify faithfulness and stability as important explanation desiderata and examine robustness under changes in input [4].

The TrustFin implementation translates several such concerns into local checks. Its faithfulness score masks selected attribution features and compares changes in model probability against a zero-vector baseline. Stability is measured as a rank correlation between SHAP contributions before and after seeded Gaussian noise in transformed feature space. Consistency is the Jaccard similarity between top-three SHAP and LIME feature sets. These are project-specific operational definitions. Their scores should not be treated as interchangeable with every published fidelity, sensitivity, or agreement metric.

\section{Trustworthy and Fair AI}
Fairness can be defined in different ways, and metrics may express different policy choices. Hardt, Price, and Srebro formalize equality of opportunity in supervised classification in terms of true-positive rates across groups [10]. TrustFin AI’s live request metric is different: it compares the model’s approval probability for one applicant with the probability after flipping the gender field while holding the remaining normalized profile fixed. That individual counterfactual score is not a population-level disparate-impact or equal-opportunity audit.

The saved training metrics also contain a female-to-male approval-rate ratio calculated on each generated test split. This is a group-level descriptive statistic for the synthetic evaluation data, while the online counterfactual score is per applicant. Neither number on its own establishes that a deployed credit system is fair. The NIST AI Risk Management Framework describes trustworthy-AI characteristics and risk-management activities as context-dependent and lifecycle-oriented [5]; TrustFin AI is a academic prototype implementation and has not undergone an institutional AI risk-management assessment.

Counterfactual explanations are another related area. Wachter, Mittelstadt, and Russell discuss explanations framed as changes to facts that could produce a desired result, while noting technical and normative constraints [11]. TrustFin’s what-if calculator does not implement a counterfactual search algorithm. It adjusts three financial inputs and recalculates an illustrative affordability ratio; it does not query the classifier or guarantee an outcome.

\section{Comparative Study of Existing Research}
Table 2.1 relates the literature to the project without attributing the prototype's custom metrics to the cited authors.

\begin{longtable}{p{0.307\linewidth}p{0.307\linewidth}p{0.307\linewidth}}
\caption{Literature and its relationship to TrustFin AI}\\

Work & Contribution & Relationship to TrustFin AI\\
Hand and Henley [1] & Credit-classification review and domain issues & Establishes the credit-scoring context; does not validate the project datasets.\\
Lessmann et al. [2] & Empirical classifier benchmarking & Motivates comparative evaluation; the product scripts use a fixed XGBoost configuration.\\
Friedman [6]; Chen and Guestrin [7] & Gradient boosting and XGBoost & Supports the model family and implementation choice.\\
Lundberg and Lee [8]; Lundberg et al. [12] & Additive attribution and tree explanations & Supports local SHAP contributions and TreeExplainer.\\
Ribeiro et al. [9] & Local surrogate explanation & Supports the LIME path and attention to neighborhood assumptions.\\
Yeh et al. [3]; Alvarez-Melis and Jaakkola [4] & Explanation infidelity, sensitivity and robustness & Motivates separate explanation checks; does not establish validity of the project's score rescaling.\\
Hardt et al. [10] & Equality of opportunity & Distinguishes a group criterion from online gender-flip sensitivity.\\
Wachter et al. [11] & Counterfactual explanation & Informs future model-verified recourse; the current calculator estimates affordability.\\
\end{longtable}
The integration contribution is the joint presentation of several explicitly defined diagnostics. Agreement between two explainers is useful review information, but cannot independently establish financial validity or causal fairness.

\section{Research Gap}
The literature provides explanation methods and several proposals for evaluating explanations, but no single scalar is universally accepted as a complete measure of explanation trustworthiness. The project addresses an engineering gap: a working loan-assessment system can show two local explanations and attach several explicit checks to each decision. The contribution is the integration and examination of these checks in a prototype system, not a new theoretical metric or evidence of real-world lending validity.

\section{Summary}
SHAP and LIME provide complementary local explanation mechanisms, while explanationevaluation research motivates checking behavior, stability, and agreement separately. Fairness metrics require a clear definition and a relevant population. These findings shape the report’s interpretation: TrustFin AI produces observable model outputs and project-defined validation scores, but the scores remain limited to the definitions, data, and perturbations implemented in the repository.

\chapter{Problem Statement and Requirements}
\section{Problem Definition}
The project addresses the difficulty of reviewing a loan model output when the classifier is more complex than a short rule list. A bare approval or rejection label gives little indication of which model inputs were influential, and one post-hoc explanation does not show whether the explanation is robust or consistent with another method. The prototype system produces an individual loan prediction, generates two explanations, computes defined explanation checks, and presents the resulting information through a web interface.

\section{Existing System}
Conventional credit scoring commonly uses structured application and repayment features with statistical or machine-learning classifiers. A model may return a class and probability, but an applicant-facing or officer-facing workflow also needs data validation, understandable reasons, persistence, and a way to inspect prior evaluations. Post-hoc explainers can contribute to this workflow, but they should not be assumed valid without a stated evaluation procedure.

\section{Limitations of a Prediction-Only Approach}
A prediction-only interface omits important context: whether input values were normalized, which features influenced the local result, whether multiple explanation methods agree, and how stable an explanation is after a defined perturbation. Likewise, presenting one scalar called confidence can be misleading if users interpret it as probability calibration or as trust in the explanation. The system therefore returns separate prediction and explanation-quality fields and labels the composite score as a project-defined indicator.

\section{Implemented System}
TrustFin AI provides two loan products, each with its own XGBoost model, preprocessor, engineered financial features, and explanation background sample. A FastAPI route normalizes an application, calculates monthly EMI and FOIR, gets a model prediction, generates SHAP and LIME explanations, evaluates the four configured checks, creates a dossier identifier, and attempts to persist an evaluation record. The React client presents an applicant view and a bank-officer view. Saved records can be inspected in history, batch CSV records can be processed, and a ReportLab service generates a PDF assessment report.

\section{Advantages}
The system makes the prediction and explanation stages inspectable in one workflow. Product-specific fields prevent property-collateral features from being treated as personalloan fields and keep revolving-credit and loan-purpose fields specific to unsecured lending. Multiple explanation checks expose more information than a single attribution chart. Local SQLite support allows use without a hosted PostgreSQL service. The record schema preserves input JSON and model outputs for later inspection.

These advantages are bounded by the generated datasets, one stored split per model, local metric definitions, and the absence of server-side authentication. TrustFin AI is a academic project prototype; its synthetic-data evaluation and access-control design define its intended operating scope.

\section{Functional Requirements}
Table 3.1 distinguishes implemented functions from broader design capabilities. A requirement describes scope; runtime evidence is reported separately in Chapter 6.

\begin{longtable}{p{0.307\linewidth}p{0.307\linewidth}p{0.307\linewidth}}
\caption{Functional requirements and prototype realization}\\

ID & Requirement & Prototype realization\\
FR-01 & Select a loan product and presentation mode & Home/personal product and applicant/officer selection.\\
FR-02 & Check and normalize input & Pydantic types/required fields plus registry defaults; financial-domain checks remain partial.\\
FR-03 & Produce classification and probability & Product XGBoost model and predict\_proba.\\
FR-04 & Explain the local prediction & TreeSHAP and LIME contributions in StandardExplanation.\\
FR-05 & Calculate explanation diagnostics & Faithfulness, stability, top-three consistency and gender sensitivity.\\
FR-06 & Summarize diagnostics separately from probability & Weighted Trust Score and predicted-class probability.\\
FR-07 & Assess repayment burden & EMI, FOIR, solvency band and officer recommendation; no backend class override.\\
FR-08 & Store and inspect evaluations & SQLAlchemy records, list/detail history and batch linkage; save failures are reported operationally rather than blocking inference.\\
FR-09 & Process application batches & CSV/JSON row validation, valid results and flawed-row reasons.\\
FR-10 & Produce a report & On-demand ReportLab PDF.\\
FR-11 & Explore financial scenarios & Local deterministic what-if calculator.\\
FR-12 & Manage authenticated users and administrative operations & Conceptual design capability; outside the current account-free prototype.\\
\end{longtable}
The authenticated and distributed capabilities depicted in some original diagrams express the proposed design envelope. They are not included among the verified functional outcomes.

\section{Non-Functional Requirements}
The implementation aims for a modular separation between React presentation, FastAPI routes, model registry, explanation components, validation functions, database access, and PDF reporting. Pydantic schemas provide request and response structure. Reproducibility is supported through a SHA-256-derived applicant seed for LIME sampling and a seeded Gaussian perturbation for the stability check. Local SQLite is configured as a zero-configuration database path, with optional PostgreSQL connectivity. These properties are implemented design choices; no service-level performance, availability, or production security target is measured in the repository.

\section{Hardware Requirements}
The prototype requires a workstation capable of running Python scientific packages and the Node.js development server. Inference and persistence checks were performed locally, with temporary database isolation. The processor, RAM and operating-system build were not captured in the retained verification evidence; no minimum capacity or concurrent-user limit is claimed. The actual development-machine specification belongs in the group's final submission record once confirmed.

\section{Software Requirements}
The backend uses Python, FastAPI, Pydantic, SQLAlchemy, pandas, NumPy, scikit-learn, XGBoost, SHAP, LIME, SciPy and ReportLab. React, Vite, Tailwind CSS, Axios, Recharts and Lucide support the interface. pytest and Vitest with Testing Library provide software tests.

The repository README lists Python 3.10+ and Node.js 18+ as general prerequisites; the retained successful local checks used Python 3.12.7 and Node.js 22.19.0. These observed versions are evidence of one working environment, not a tested compatibility matrix. Most Python dependencies have lower bounds. Loading preprocessors saved with scikit-learn 1.7.2 under installed 1.9.1 produced version warnings. Appendix A records the observed versions.

\section{Project Constraints}
The main constraints are the synthetic nature of the two product datasets, the limited hold-out evaluation evidence, the dependence of explanations on preprocessing and perturbation assumptions, and the lack of an authenticated access-control layer. The Trust Score is weighted according to one project configuration and has no external calibration study. Population fairness, causal validity, and live lender integration are outside the implemented scope.

The current financial-domain input checks and advisory-only solvency handling constrain interpretation of results. The available verification establishes selected local operations, not restart durability, PostgreSQL failover under outage, or an interactive browser walkthrough.

\chapter{System Analysis and Design}
\section{System Overview and Design Boundary}
TrustFin AI is organized around presentation, application orchestration, prediction, explanation, trust evaluation and persistence. The original diagrams express the proposed design at different abstraction levels. The local prototype realizes the central prediction-explanation-validation path in React and FastAPI, with product artifacts and SQLAlchemy records.

The diagrams are reproduced without editing their labels, arrows, colours or layout. A single design-to-prototype mapping in Table 4.1 establishes how to read them. Terms such as controller, evaluator and report generator can denote responsibilities even when implemented as module-level functions rather than classes.

\begin{longtable}{p{0.46\linewidth}p{0.46\linewidth}}
\caption{Conceptual design to prototype mapping}\\

Conceptual element & Local prototype realization\\
Dashboard / frontend & App.jsx and product/role-specific React components.\\
PredictionController / application service & prediction.py route functions and evaluate\_single\_record().\\
MLModel / prediction component & ModelRegistry, ProductPipeline and saved XGBoost artifacts.\\
Explainability & SHAPExplainer, LIMEExplainer and StandardExplanation.\\
TrustEvaluator / trust engine & Four validation functions plus calculate\_trust\_score().\\
ReportGenerator & generate\_loan\_report\_pdf() in the reporting module.\\
Data management & PredictionLog and BatchUpload through SQLAlchemy sessions.\\
Authentication / administrative management & Proposed design responsibilities, not implemented accounts or protected APIs.\\
Separate AI server and HTTPS & Proposed deployment topology; local execution uses in-process ML and HTTP development servers.\\
\end{longtable}
Where a diagram records an earlier evaluation configuration, the implementation formulas in Chapter 5 and results in Chapter 7 define the tested prototype. This preserves design history while keeping the experiment reproducible.

\section{Conceptual System Architecture}
Figure 4.1 organizes the framework into a presentation layer, API orchestration, decision support, dual explanation, validation and reporting. An application enters through the frontend and its financial features are transformed before model inference. SHAP and LIME produce local contributions, and validation values accompany the returned assessment.

\begin{figure}[p]\centering\includegraphics[width=\linewidth,height=.82\textheight,keepaspectratio]{\detokenize{assets/system architecture.png}}\caption{Original conceptual system architecture of TrustFin AI.}\end{figure}
The prototype uses the same central responsibilities, with two product pipelines selected by the model registry. Several labels preserve an earlier design configuration: the image shows equal 25\% weighting, equal-opportunity fairness, top-three faithfulness, a calibrated model and an override check. Current execution instead uses 35/25/30/10 weights, individual gender-flip sensitivity, threshold-selected SHAP masking, uncalibrated class probabilities and separate solvency metadata. Database writes are ordinary route/session operations rather than an asynchronous logging service. These differences locate the prototype within the proposed architecture without changing the original image.

\section{Methodology and Overall Workflow}
The methodology begins with loan records, preprocessing and model construction, then proceeds through inference, explanations, trust evaluation and decision-support presentation. Figure 4.2 retains the original high-level methodology.

\begin{figure}[p]\centering\includegraphics[width=\linewidth,height=.82\textheight,keepaspectratio]{\detokenize{assets/methodology_render.png}}\caption{Original proposed methodology from loan data to explanation-supported decision review.}\end{figure}
In the product experiments, the record source is synthetic generation rather than institutional historical lending data. Training uses a fixed model configuration; the tuning box denotes a broader design activity rather than an executed parameter-search experiment. Online execution starts from application input and reuses saved models and preprocessing artifacts. It does not retrain during prediction.

\section{Evaluation Flow and Data Flow Analysis}
For a single application, the API checks request structure, selects the product pipeline, normalizes fields, derives financial features and transforms one record. It then obtains the model class and probability, computes both explanations, calculates four diagnostics and assembles the response. A persistence attempt follows computation. Invalid required fields can produce HTTP 422; successful type checking alone does not establish a financially valid domain.

At context level, the applicant or officer supplies application information and receives an assessment or report. At a more detailed level, the data move through input processing, product transformation, inference, explanation, diagnostic evaluation and optional storage. History requests retrieve stored JSON and scalar summaries; PDF requests carry input and result data to the reporting function. Batch input is a collection of rows, each classified as usable or flawed before evaluation.

The supplied asset set has no separate context DFD, Level 1 DFD, activity diagram or detailed flowchart. Their analytical content is therefore explained here without substituting an unrelated image or inventing an original diagram. The original workflow and sequence figures provide the available visual views.

\section{Use Case Design}
Figure 4.3 captures the proposed interactions between a user, administrator and TrustFin functions. Its prediction, result, explanation and download use cases correspond to the central prototype workflow.

\begin{figure}[p]\centering\includegraphics[width=\linewidth,height=.82\textheight,keepaspectratio]{\detokenize{assets/Use_Case_Updated.png}}\caption{Original proposed use-case diagram, including wider administrative capabilities.}\end{figure}
The working interface implements applicant and officer presentation modes rather than authenticated user/admin identities. Registration, login/logout, user management, dataset administration and retraining controls shown in the design are not exposed as implemented account workflows. The officer view offers audit/history and batch access; the applicant view offers assessment, submitted values, report actions and a financial calculator for rejected cases. Detailed trust pillars and dual-explainer charts belong to the officer audit view; PDF output also includes diagnostic values.

\section{Class and Domain Design}
The class views describe conceptual application, prediction, explanation, trust and reporting responsibilities. Figure 4.4 is the supplied final domain-oriented view; Figure 4.5 preserves the second supplied class-design view.

\begin{figure}[p]\centering\includegraphics[width=\linewidth,height=.82\textheight,keepaspectratio]{\detokenize{assets/Class_Diagram_Final.png}}\caption{Original conceptual class design: application, prediction, explanation and trust.}\end{figure}
\begin{figure}[p]\centering\includegraphics[width=\linewidth,height=.82\textheight,keepaspectratio]{\detokenize{assets/class diagram.png}}\caption{Original service-oriented class-design variant.}\end{figure}
The implemented Python classes are LoanApplicationRequest, LoanPredictionResponse, FeatureContribution, StandardExplanation, TrustEvaluation, ModelRegistry, ProductPipeline, SHAPExplainer, LIMEExplainer, PredictionLog and BatchUpload. Conceptual User and LoanApplication entities in the figures should not be read as database tables: no user table is implemented. Similarly, controller/service/evaluator responsibilities are often functions. This mapping is preferable to claiming that the source declares every conceptual class verbatim.

\section{Sequence and Activity Design}
Figure 4.6 shows application submission, input processing, prediction, SHAP and LIME calls, trust evaluation and an optional report request. This order broadly matches the sequential explanation generation in the route.

\begin{figure}[p]\centering\includegraphics[width=\linewidth,height=.82\textheight,keepaspectratio]{\detokenize{assets/Sequence_Diagram_Updated.png}}\caption{Original single-application sequence with optional report generation.}\end{figure}
In the prototype, product routing and financial calculations occur inside evaluate\_single\_record(). The response contains the raw classifier decision together with the solvency result. Persistence is attempted by the route before returning the result. Report generation is a separate API request. The logical activity path is portal selection, product selection, form entry, evaluation and role-specific review, followed optionally by history inspection or PDF export. Batch processing is an officer-facing alternative intake path.

\section{Component Structure}
Figure 4.7 groups the frontend, API, prediction, explanation and trust responsibilities. The decomposition is useful even though the prototype runs these Python responsibilities within one backend process.

\begin{figure}[p]\centering\includegraphics[width=\linewidth,height=.82\textheight,keepaspectratio]{\detokenize{assets/Component_Diagram.jpeg}}\caption{Original high-level software component design.}\end{figure}
The database symbol illustrates the proposed shared persistence boundary. Actual explainers and validation functions do not independently query PostgreSQL: route functions coordinate persistence through SQLAlchemy. Local SQLite is supported alongside optional PostgreSQL. Reporting and product routing are additional concrete modules described in Chapter 5.

\section{Deployment Design}
Figure 4.8 presents a proposed distributed deployment with a client, application server, AI server and database server.

\begin{figure}[p]\centering\includegraphics[width=\linewidth,height=.82\textheight,keepaspectratio]{\detokenize{assets/deployment-diagram.png}}\caption{Original proposed distributed deployment topology.}\end{figure}
The tested local prototype runs Vite on port 5173 and FastAPI on port 8000 over HTTP. Model inference and explanation generation occur in the FastAPI Python process, using repository artifacts rather than an independent AI REST service. SQLAlchemy supports local SQLite or a configured PostgreSQL URL. HTTPS termination and distributed hosting in the conceptual image are architectural options, not verified deployment outcomes.

\section{Package Structure}
Figure 4.9 provides a logical package view spanning presentation, application, prediction, explanation, trust and data management.

\begin{figure}[p]\centering\includegraphics[width=\linewidth,height=.82\textheight,keepaspectratio]{\detokenize{assets/Package_Diagram.jpeg}}\caption{Original proposed package decomposition.}\end{figure}
The prototype maps these responsibilities to frontend/src and backend/app subdirectories. The actual data layer consists of db/models.py and db/session.py rather than separate user, loan and explanation repository classes. Authentication packages remain part of the broader proposed design. The package view therefore communicates separation of responsibilities without requiring a one-to-one source package for each label.

\section{Database Design}
The implemented relational store has two principal entities. JSON values preserve detailed records, while scalar fields support retrieval and summary display. Table 4.2 states the actual association instead of implying a foreign-key constraint that is not declared.

\begin{longtable}{p{0.307\linewidth}p{0.307\linewidth}p{0.307\linewidth}}
\caption{Implemented persistence entities}\\

Entity & Key fields and stored content & Relationship\\
PredictionLog / prediction\_logs & Integer id; unique dossier\_id; timestamp; product; financial summaries; normalized input JSON; prediction/probability/confidence; solvency; SHAP/LIME JSON; four scores; trust; full\_response; optional batch\_id. & Optional batch\_id associates the record with a batch.\\
BatchUpload / batch\_uploads & Integer id; unique batch\_id; product; filename; timestamp; total/valid/flawed counts; flawed\_data JSON; summary JSON. & One batch can be located alongside many logs by matching the string identifier.\\
\end{longtable}
There is no declared foreign key, ORM relationship, User table or officer-signature entity. init\_db() creates tables and performs additive checks for selected columns; no versioned migration framework is used. The application-level association is adequate for prototype retrieval but is not a claim of enforced relational cascade integrity.

\section{API and Module Design}
The loan router is mounted under /api/v1/loan. Table 4.3 lists implemented routes; paths below are relative to that prefix except /health.

\begin{longtable}{p{0.307\linewidth}p{0.307\linewidth}p{0.307\linewidth}}
\caption{Implemented API routes}\\

Method & Route & Purpose\\
POST & /predict & Single application evaluation.\\
POST & /predict-batch & Multipart CSV and loan\_type.\\
POST & /predict-batch-json & JSON rows, filename and optional loan\_type.\\
GET & /evaluations & Recent evaluation summaries.\\
GET / DELETE & /evaluations/\{evaluation\_id\} & Retrieve or delete an integer-ID record.\\
GET & /batches & Recent batch summaries.\\
GET / DELETE & /batches/\{batch\_id\} & Retrieve or delete a string-ID batch and associated logs.\\
GET & /batch-sample-csv & Product CSV sample.\\
GET & /batch-presets/\{preset\_id\} & home\_loan or personal\_loan demo rows.\\
POST & /generate-pdf & PDF from supplied applicantData and predictionResult.\\
GET & /health (root) & Fixed application health response.\\
\end{longtable}
Pydantic schemas define the response contract. Missing required fields yield 422; absent records yield 404; execution errors can yield 500. Health is a liveness response rather than a full model/database readiness test. The API has no authenticated ownership boundary in this prototype.

\chapter{System Implementation and Methodology}
\section{Technology Stack and Development Environment}
Python/FastAPI/Pydantic implement the API; SQLAlchemy handles persistence; pandas and NumPy handle data; scikit-learn transforms features; XGBoost classifies records; SHAP and LIME explain them; SciPy computes rank correlation; ReportLab generates PDFs. React/Vite, Tailwind CSS, Recharts, Lucide and Axios provide presentation and communication.

The observed environment was Python 3.12.7, Node.js 22.19.0 and npm 10.9.3. Appendix A records retained library versions. No dependencies were installed or upgraded for verification. Loading preprocessors saved under scikit-learn 1.7.2 with installed 1.9.1 produced version warnings. Success on selected checks is not a general serialization guarantee.

The API entry point is backend/app/main.py; orchestration is in backend/app/api/routes/prediction.py. Relative artifact paths make the working directory significant. The frontend service hardcodes localhost:8000 rather than using the example VITE\_API\_BASE\_URL.

\section{Dataset Preparation and Provenance}
scripts/generate\_homeloan\_dataset.py and scripts/generate\_personalloan\_dataset.py use seed 42 and generate 3,000 records each. Home loans have 18 total columns and personal loans 17, including ID and target. No empty cells or duplicate full rows occur in the checked product CSVs. Loan\_ID is excluded from prediction.

The generators sample demographics and financial variables from specified distributions. Income depends partly on education and is clipped to product ranges; existing EMI is a sampled fraction of income. CIBIL and missed payments are correlated. Home-loan amounts/terms relate to income, age and collateral. Personal loans include purpose and utilization effects.

Labels combine a latent score with explicit rules. Home-loan hard rejection includes CIBIL below 560, at least three missed payments, FOIR above 0.68 or LTV above 92\%; a strong-profile rule can approve specified low-burden records. Personal-loan rejection includes CIBIL below 580, at least three missed payments, FOIR above 0.60, utilization above 0.88 or income below 19,000. These generator thresholds differ from the online 50\% solvency classification.

Both generators subtract 0.65 from the latent score of female applicants in a borderline band before assigning labels. This deliberately creates demographic sensitivity in the simulation. It is not evidence about an institution. datasets/train.csv and data/processed/train\_clean.csv belong to the legacy path; original external provenance of train.csv is unresolved and is not used to support the product experiments.

\section{Dataset Fields and Feature Representation}
Table 5.1 groups raw inputs. Income and EMI are monthly rupee amounts; principal and property value are rupee totals. Ratios and percentages are distinguished.

\begin{longtable}{p{0.307\linewidth}p{0.307\linewidth}p{0.307\linewidth}}
\caption{Product input field groups}\\

Field group & Home loan & Personal loan\\
Demographics & Gender, Married, Dependents, Education, Age & Same\\
Employment & Employment\_Status, Work\_Experience & Same\\
Income/debt & Applicant\_Income, Coapplicant\_Income, Existing\_EMI & Applicant\_Income, Existing\_EMI\\
Facility & Loan\_Amount, Loan\_Term in months & Same\\
Bureau & CIBIL\_Score, Missed\_Payments\_12M & Same\\
Product-specific & Property\_Value, Property\_Area & Loan\_Purpose, Credit\_Card\_Utilization\\
Excluded & Loan\_ID, Loan\_Status & Loan\_ID, Loan\_Status\\
\end{longtable}
Categorical columns are Gender, Married, Education, Employment\_Status and Property\_Area for home loans; personal loans substitute Loan\_Purpose for area. OneHotEncoder drops the first category and ignores unknown categories. Numerical columns use StandardScaler. Explicitly selected model vectors contain 20 home-loan and 21 personal-loan features. This is domain selection, not a learned feature-selection experiment.

\section{Preprocessing and Request Normalization}
The scripts derive financial features, exclude ID/target, map Approved to 1 and Rejected to 0, and create an 80/20 stratified split with seed 42. Preprocessing is fitted on training predictors and applied to held-out records without refitting. No product imputation pipeline is present; the checked product CSVs are complete.

LoanApplicationRequest requires Applicant\_Income, Loan\_Amount and Loan\_Term. Optional numeric validators coerce values. ModelRegistry.normalize\_input() supplies defaults, converts legacy Credit\_History to CIBIL where appropriate, estimates missing property value as Loan\_Amount/0.75 and supplies missing bureau values. Malformed optional CIBIL can become 720; malformed existing EMI can become zero.

These rules explain what reaches the model, but single-record requests do not comprehensively enforce positive values, valid ages or product tenures. Batch intake adds row-level ranges and required-field checks. Chapter 6 separates successful structural validation from observed financial-domain failures.

\section{Exploratory Description and Class Distribution}
The evidence consists of generated CSVs, generator definitions and saved evaluation artifacts rather than a separate EDA notebook. The descriptive analysis reports observed counts and construction relationships without claiming a comprehensive statistical study.

\begin{figure}[p]\centering\includegraphics[width=\linewidth,height=.82\textheight,keepaspectratio]{\detokenize{assets/class_counts.png}}\caption{Approved and rejected records in the two synthetic product datasets.}\end{figure}
Home loans contain 1,529 approved and 1,471 rejected labels; personal loans contain 1,563 approved and 1,437 rejected labels. The near-balanced classes reflect generator choices. Full-row uniqueness does not establish a complete feature-level leakage audit or external validation.

\section{Financial Feature Engineering}
For principal P, annual decimal rate i and term n months, r=i/12. The product functions calculate fixed-rate amortized installment:

\begin{equation}\mathrm{EMI}=P\,\frac{r(1+r)^n}{(1+r)^n-1}\end{equation}
Home loans use i=0.085; personal loans use i=0.135. EMI is rounded to two decimals. These are illustrative feature-engineering rates, not live quotations or validated bank policy. Income I is applicant plus co-applicant income for home loans and applicant income for personal-loan FOIR.

\begin{equation}\mathrm{FOIR}=\frac{\mathrm{Existing\ EMI}+\mathrm{Proposed\ EMI}}{I}\end{equation}
FOIR is rounded to four decimals. Home loans additionally calculate:

\begin{equation}\mathrm{LTV}(\%)=100\,\frac{\mathrm{Loan\ Amount}}{\mathrm{Property\ Value}}\end{equation}
LTV is rounded to two decimals. Total\_Income is also derived for home loans, although the model uses explicitly selected columns. Credit\_Card\_Utilization is a fraction; the request validator divides percentage-like values above 1 and at most 100 by 100.

Zero denominators and terms are replaced with safe values to prevent selected arithmetic failures; this does not make invalid inputs meaningful. Legacy EMI is Loan\_Amount/Loan\_Term without interest, with proposed-EMI-only DTI. It must not be substituted for the product calculations.

\section{Deterministic Solvency Assessment}
evaluate\_single\_record() produces the solvency summary in Table 5.2.

\begin{longtable}{p{0.307\linewidth}p{0.307\linewidth}p{0.307\linewidth}}
\caption{Project solvency bands}\\

FOIR & Status & isSolvent\\
At most 0.35 & OPTIMAL & True\\
Above 0.35 and at most 0.50 & MODERATE & True\\
Above 0.50 & INSOLVENT & False\\
\end{longtable}
The response includes proposed/existing EMI, total debt, income, ratio, formatted percentage and warning. Total monthly obligation is proposed plus existing EMI; netDisposableIncome is reported income less debt payments. Monthly living expenses are not collected, so this is not surplus after all household expenditure.

The 50\% threshold is the project's affordability guide. Some UI/PDF wording calls it regulatory; no institutional policy validation is established. The backend preserves the classifier's result even when isSolvent is false. Officer guidance can recommend declining an approved/insolvent case without changing the API class.

Personal-loan FOIR uses applicant income. A caller can supply co-applicant income that appears in reported total income while the FOIR denominator remains unchanged. The ordinary personal-loan form sends zero co-applicant income. Legacy DTI excludes existing EMI although total obligations can include it.

\section{XGBoost Model Development}
scripts/train\_homeloan.py and scripts/train\_personalloan.py train XGBClassifier with 120 estimators, depth 4, learning rate 0.08, subsample 0.85, colsample\_bytree 0.85, random\_state 42 and logloss evaluation metric [7]. This is fixed configuration, without implemented parameter search, cross-validation or probability calibration.

Algorithm 5.1: Product model training.

1. Read the product CSV and compute financial features.

2. Exclude Loan\_ID and Loan\_Status from predictors; encode target labels.

3. Split into stratified 80/20 subsets using seed 42.

4. Fit encoding/scaling on training predictors; transform both subsets.

5. Fit XGBoost with the fixed configuration.

6. Calculate held-out accuracy, precision, recall, F1 and ROC-AUC.

7. Calculate female/male predicted approval-rate ratio.

8. Save model, preprocessor, feature names, background array and metrics.

The background array is the first 100 transformed training rows. ModelRegistry loads these artifacts into ProductPipeline with the feature function and explainers. The legacy path remains for compatibility. Route version strings such as home\_loan\_v2 are identifiers rather than artifact hashes.

\section{Evaluation and Artifact Loading}
Each product has 600 held-out records. metrics.json stores accuracy, precision, recall, F1, ROC-AUC, descriptive gender approval ratio and feature count. Product confusion matrices are not saved in those files. Legacy model\_metadata.json is a separate artifact and is not substituted for product results.

The registry initializes available pipelines at import. Unknown product identifiers can fall back to home loan. Requests without explicit product use compatibility routing based partly on bureau fields; the frontend supplies a product. The route also loads legacy artifacts separately and gates endpoints on that global model, so product operation is not completely independent of legacy initialization. /health is a fixed liveness response, not a model/database readiness check.

\section{SHAP Explanation Generation}
SHAPExplainer initializes shap.TreeExplainer(model). explain\_instance() requests contributions for the transformed vector and returns StandardExplanation: method plus FeatureContribution records with name, value, contribution and direction.

For the loaded binary XGBoost classifiers, default TreeExplainer output is interpreted in raw-margin/log-odds space toward approval. Positive values support approval; negative values reduce it. Contributions between -0.0001 and 0.0001 are neutral. Rejected cases retain this approval-oriented interpretation.

The response omits the SHAP base value. Raw values are substituted only when feature names match original-values keys; one-hot coordinates remain encoded and derived LTV\_Ratio can remain standardized. UI helpers attempt further mapping. These are model attributions, not causal effects [8,12].

\section{LIME Explanation Generation}
LIMEExplainer wraps LimeTabularExplainer in classification mode using transformed background and feature names [9]. model.predict\_proba evaluates neighbors; the wrapper requests all features and extracts exp.as\_map()[1], the approved-class explanation even for rejected cases.

No explicit sample count, kernel width or categorical indices are supplied, so installed library defaults apply. Product background is the saved 100-row sample. Applicant seeding resets explainer/base/discretizer random states where supported and also sets NumPy's global seed.

LIME weights describe a local surrogate and differ in scale from SHAP log-odds. The common schema permits comparison of feature names/signs, not interchangeable magnitudes. Surrogate intercept, fit score and local prediction are not returned, so output generation alone does not establish approximation quality.

\section{Explanation Validation Framework}
The pillars have distinct targets: faithfulness/stability use SHAP, consistency compares SHAP and LIME, and fairness evaluates model gender sensitivity. LIME faithfulness/stability are not separately measured. These operational definitions are motivated by explanation-evaluation research [3,4], not presented as validated replacements for published metrics. Generation and validation execute sequentially in evaluate\_single\_record().

\subsection{Faithfulness}
calculate\_faithfulness() obtains original approval probability p, predicted class and zero-vector probability p0. For approval it selects contributions above 0.08, falling back to up to five positive contributors. For rejection it selects contributions below -0.08, falling back to up to five negative contributors by absolute magnitude. The API does not pass top\_k. Selected transformed coordinates are zeroed; pm is the masked probability.

\begin{equation}q=\frac{|p-p_m|}{\max(0.08,|p-p_0|)}\end{equation}
\begin{equation}F=\max(0.02,2q),\quad q\leq 0.05\end{equation}
\begin{equation}F=0.74+0.17(1-e^{-q}),\quad q>0.05\end{equation}
The result is capped at 0.95 and rounded to four decimals; no salient features yields 0.10. This is heuristic rescaling, not calibrated explanation correctness. Absolute change does not verify expected direction. The mapping jumps sharply above q=0.05. Zero means a scaled numeric mean or baseline category, and masking derived coordinates can break financial relationships. The zero vector is not necessarily a realistic applicant baseline.

\subsection{Stability}
An input-derived seed controls one Gaussian perturbation with standard deviation 0.01 across all transformed coordinates, including one-hot variables. Original and perturbed SHAP vectors are compared using Spearman rank correlation rho:

\begin{equation}S=\frac{\rho+1}{2}\end{equation}
Undefined correlation returns zero. This is not a 1\% change to each raw financial field. It need not preserve categorical validity or financial relationships. Rank stability also does not establish magnitude stability or robustness throughout a neighborhood.

\subsection{Consistency}
Both explanations are ranked by absolute contribution. A3 and B3 contain their three leading feature names:

\begin{equation}C=\frac{|A_3\cap B_3|}{|A_3\cup B_3|}\end{equation}
Overlap counts 0, 1, 2 and 3 yield 0, 0.2, 0.5 and 1. The score ignores signs, magnitudes and ordering within the sets. The officer audit's directional-agreement display is a separate frontend statistic.

\subsection{Individual Gender-Flip Fairness Score}
calculate\_counterfactual\_fairness() scores normalized x, flips Male/Female, recomputes financial features and preprocessing, and scores xg with other normalized fields fixed:

\begin{equation}R=\max(0,\min(1,1-|p(x)-p(x^{(g)})|))\end{equation}
This is binary gender sensitivity, not causal fairness, group parity, equal opportunity or compliance. A small change can cross a decision boundary while still yielding a high score. Exceptions return 1.0; an executed error-path check confirmed that current fallback behavior. The equal-opportunity helper is separate from the live request pillar [10].

\section{Composite Trust Score and Prediction Probability}
Default weights in calculate\_trust\_score() are:

\begin{equation}T=0.35F+0.25S+0.30C+0.10R\end{equation}
The function does not normalize custom weights. UI/storage bands use at least 0.85 for high trust, at least 0.70 for moderate trust and lower values for low trust. T informs frontend recommendations but does not change backend prediction. No externally validated threshold study is claimed.

For approval probability p and predicted class yhat, confidenceScore is:

\begin{equation}Q=p\ \mathrm{if}\ \hat{y}=1;\quad Q=1-p\ \mathrm{otherwise}\end{equation}
Q is predicted-class probability, not T and not verified calibration. Neither isotonic regression nor Platt scaling is implemented in product training. Faithfulness rescaling does not calibrate the classifier.

\section{Backend Orchestration and API Flow}
FastAPI mounts the loan router under /api/v1/loan. predict\_loan() parses LoanApplicationRequest and calls evaluate\_single\_record(). Algorithm 5.2 summarizes its path.

1. Select ProductPipeline and normalize application fields.

2. Derive financial features and transform the record.

3. Obtain model class and approval probability.

4. Calculate solvency metadata separately from the class.

5. Generate SHAP and approval-class LIME explanations.

6. Calculate masking, seeded stability, top-three overlap and gender-flip scores.

7. Combine T and calculate predicted-class probability Q.

8. Create dossier metadata and database fields.

9. Attempt persistence and return the structured result.

The route adds/commits/refreshes PredictionLog and assigns evaluationId. Save errors are caught and rollback attempted, so a successful prediction can lack persistence. Evaluation errors can expose a 500 detail string; the custom 422 handler logs invalid values. An async route declaration does not make CPU-heavy explainers concurrent or create a background job system.

\section{Frontend Components and User Workflows}
App.jsx manages LANDING, LOAN\_SELECTION, APPLICATION\_FORM and RESULTS states. LandingPage selects mode; LoanTypeSelection selects product; ApplicantForm collects fields or loads samples and calls predictLoan(). Loading indicators and error states cover prediction, history, batch and PDF requests. Form submission substitutes defaults for several empty/zero values.

\begin{longtable}{p{0.46\linewidth}p{0.46\linewidth}}
\caption{Active frontend responsibilities}\\

Component & Responsibility\\
BankOfficerSummary & Decision, affordability, leading factors and recommendation.\\
BankOfficerAudit & SHAP/LIME charts and agreement comparison.\\
ApplicantPortalView & Assessment, submitted form and PDF actions.\\
RecourseSimulator & Rejected-case financial sliders.\\
AuditHistoryModal & Product-filtered evaluations and batch history.\\
BatchUploadModal & CSV/demo intake, valid and flawed rows.\\
CustomerNoticeModal & Officer customer-notice preview.\\
SubmittedApplicationFormModal & Submitted values.\\
services/api.js & Axios prediction, history, batch and PDF calls.\\
\end{longtable}
Full diagnostic charts are officer-facing. The applicant screen uses selected adverse-factor conditions rather than rendering trust pillars. InsightsDashboard.jsx is retained but not imported by the active shell. Roles are client state. Some value helpers retain legacy CIBIL/EMI assumptions, so input/response records are used for the worked examples.

\section{Database Persistence and Application History}
PredictionLog stores normalized input JSON, financial summaries, classifier results, SHAP/LIME objects, scores, trust/rating and full\_response. BatchUpload stores metadata, counts, flawed rows and summary. Timestamps are serialized as explicit UTC for local rendering. Dossier/batch references combine product prefixes with UUID fragments.

History retrieves stored responses without explanation recomputation where full\_response exists. evaluationId is assigned after response preparation and restored during detail retrieval. Legacy records have fallback reconstruction. Default list limits provide recent records, not a complete institutional archive.

session.py loads DATABASE\_URL and creates a SQLite fallback. Startup init\_db() can create tables and add selected columns; get\_db() falls back after primary connection failure. These are operational alternatives, not replicated databases. No foreign key is declared for batch\_id. Delete APIs exist, so the ledger is not immutable. Verification used in-memory engines, leaving the existing database untouched.

\section{Batch Processing}
/predict-batch parses UTF-8 CSV through DictReader; /predict-batch-json accepts rows directly. validate\_batch\_record() normalizes aliases, infers/accepts product, checks required values/ranges and supplies product defaults. Flawed rows retain row index, raw values and reason.

process\_rows\_into\_batch() sequentially evaluates usable rows through the full pipeline, flushes prediction logs, accumulates counts/mean trust and commits a batch record. This is not vectorized inference or a queued job system. Evaluation errors may become flawed-row reasons; commit failures are caught. Returned results therefore do not alone guarantee storage.

Current product presets contain nine rows including one missing-income case. Sample CSV output appends a tenth row. Old tests expect batch\_1/batch\_2 and different counts; Chapter 6 separates these stale expectations from successful current product execution.

\section{Recommendations and What-If Calculation}
getUnderwritingAction() in decisionAdvisory.js returns high-risk/decline guidance for insolvency, or for rejection with predicted-class confidence at least 65\%. Approval with confidence at least 70\%, trust at least 70\%, solvency and FOIR at most 35\% gives a sanction recommendation; other cases suggest officer review. These labels do not persist an officer decision or alter classifier output.

RecourseSimulator recomputes EMI, FOIR, surplus and home-loan LTV. Home controls vary amount, term and property value; personal controls vary amount, term and existing EMI. Fixed illustrative rates apply. Slider changes do not call predictLoan(). The calculator supports affordability comparison, not minimal-change optimization or guaranteed model approval [11].

\section{Authentication and Security Boundary}
The prototype has no registration/login, passwords, tokens or server-enforced role access. Applicant/officer selection changes visible controls. CORS allows all origins and credentials. Authentication/admin boxes in Chapter 4 express intended design responsibilities rather than completed controls.

Financial JSON is stored and history/deletion APIs lack ownership checks. Operational deployment would require access, privacy, retention and security controls. No live bureau integration, Bank of Baroda integration or official policy validation is established.

\section{PDF Report Generation}
generate\_loan\_report\_pdf() builds an A4 ReportLab document in BytesIO with reference, applicant/facility summary, assessment, financial values, pillar weights/scores, guidance and disclaimer. Segoe UI supports rupee glyphs on Windows where available; fallback fonts are configured.

/generate-pdf accepts applicantData and predictionResult from the client rather than retrieving a server-authoritative determination. Blob URLs support preview/download. Signature/seal elements are presentation placeholders, not a digital signature; SHA-256 elsewhere derives random seeds, not a PDF audit hash.

Some narratives are static: approved/insolvent results can select sanction wording, living-expense surplus is not based on expense inputs, and product next steps are not uniformly tailored. This is a provisional computer-generated assessment, not an executed lending sanction.

\chapter{Testing and Validation}
\section{Verification Method and Status}
Phase 1 verification on 4 November 2026 combined source inspection, artifact evaluation, selected automated tests, HTTP checks and sample PDF inspection. Database writes were isolated in memory and Vite caches placed temporarily. Source, existing database and dependencies were preserved.

VERIFIED means successful execution of the stated case; PARTIALLY VERIFIED identifies a narrower boundary; NOT WORKING records failure; NOT TESTED denotes unexecuted behavior; SOURCE-CODE CONFIRMED denotes inspected implementation. Model accuracy, test-file presence and system-test outcomes are separate evidence categories.

\section{Backend Automated Tests}
Selected pytest execution covered test\_audit\_reproducibility.py, test\_data\_and\_solvency.py, test\_multiproduct\_api.py, test\_reporting\_and\_pdf.py, test\_shortcut\_learning.py and test\_xai\_and\_ablation.py: 14 passed, 7 warnings. This is not the full backend suite. Engines/sessions were redirected in-process to SQLite memory before application use.

Five additional batch functions were invoked individually: one passed and four failed.

\begin{longtable}{p{0.307\linewidth}p{0.307\linewidth}p{0.307\linewidth}}
\caption{Additional existing batch checks}\\

Check & Outcome & Observation\\
Valid legacy-style row & Failed & Expects cleaned Credit\_History, absent from current product cleaning.\\
Missing income & Passed & Row rejected with required-income reason.\\
Missing credit history & Failed & Old expectation differs from current CIBIL handling.\\
Sample CSV count & Failed & Old count differs from product sample.\\
Old preset identifiers & Failed & batch\_1/batch\_2 differ from home\_loan/personal\_loan.\\
\end{longtable}
Passing assertions are interpreted narrowly. Rejection of the extreme-debt sample does not prove override because the model itself rejects it. The ablation test supplies hardcoded degraded metrics to the trust function; it does not retrain or reproduce printed historical accuracy. Warnings included serialization-version mismatch and TestClient deprecations.

\section{Frontend Automated Tests}
Vitest ran four files covering API service, applicant form, history and batch components: 18 passed and 7 failed; two files passed and two failed. API-service and applicant-form files passed. History/batch failures include fixtures lacking product fields required for current filtering and expectations for old demo labels.

These are mocked/jsdom tests, not browser/backend end-to-end tests. Failed expectations are retained as failures. They reveal regression-alignment work without establishing that every underlying live workflow is nonfunctional.

\section{API, Batch and Persistence Checks}
Table 6.2 records actual checks against the isolated backend.

\begin{longtable}{p{0.307\linewidth}p{0.307\linewidth}p{0.307\linewidth}}
\caption{API and persistence verification}\\

Check & Status & Observation\\
Startup/health/OpenAPI & VERIFIED & HTTP 200 and routes available.\\
Both product predictions & VERIFIED & Class, explanations, solvency and pillars returned.\\
Identical HTTP payload twice & VERIFIED for sequence & Trust and LIME objects matched.\\
Save/list/detail & PARTIALLY VERIFIED & In-memory success, not restart durability.\\
Product JSON presets & VERIFIED & Each: 9 rows, 8 valid, 1 flawed.\\
CSV sample & VERIFIED & 2 rows: 1 valid, 1 flawed.\\
Omitted required fields & VERIFIED & HTTP 422.\\
Negative income/tenure & NOT WORKING as domain validation & HTTP 200, negative EMI and OPTIMAL status possible.\\
Malformed optional values & Limited validation & Defaults replace malformed values.\\
Approved above 50\% FOIR & VERIFIED behavior & Approved with INSOLVENT status.\\
Fairness exception & VERIFIED failure mode & Returned 1.0 after error.\\
\end{longtable}
No deletion was executed. Persistent-database writes, PostgreSQL connectivity, outage failover and restart durability were not tested. Batch detail retrieval succeeded in memory.

\section{Model and Explanation Checks}
Saved product artifacts were evaluated on the deterministic script-defined split without retraining. Recomputed accuracy/ROC-AUC matched stored values: home 0.97/0.9978102352; personal 0.9733333333/0.9963375672. Other classification metrics are identified as saved values.

The pipeline produced 20 home and 21 personal SHAP/LIME features and all four diagnostics. Identical HTTP payloads matched in tested sequential execution. Direct dictionaries and schema-expanded requests can hash differently because defaults alter serialized input. Reproducibility refers to identical computational input in a fixed environment, not every semantically equivalent payload.

Successful computation does not validate the heuristic score scale, causal fairness or aggregate explanation quality.

\section{Manual and PDF Verification}
Vite returned HTML and transformed React modules with HTTP 200. The browser tool failed twice before opening the application. Interactive walkthrough, chart interaction and real browser-to-backend verification therefore remain untested.

Four PDF samples covered approval/rejection for both products. All returned HTTP 200, application/pdf and one A4 page around 45 KB. Text parsing checked content and a rendered home-loan approval sample was legible. This is sample testing, not exhaustive layout verification. Chapter 7 includes a clearly labelled application report specimen from the retained evaluation, not an invented browser screenshot.

\section{Verification Summary}
The central computational path operates in the tested environment. Evidence includes 14 selected backend passes, 1 pass/4 failures among extra batch checks and 18 frontend passes/7 failures. Missing required fields are rejected, while financial-domain validation is partial. Solvency is advisory. Appendix B maps evidence to source modules.

Counts are retained from Phase 1 rather than represented as a fresh full-suite run. Raw terminal logs were not retained as independent files; this draft records their provenance and scope. No claim of production deployment, user authentication or external lending performance follows from these checks.

\chapter{Results and Discussion}
\section{Experimental Setup and Evidence Provenance}
Each product experiment uses 3,000 generated records, an 80/20 stratified split and seed 42. Stored metrics originate from the training scripts; Phase 1 recomputed accuracy/ROC-AUC using saved artifacts without retraining. Tables distinguish this evidence from selected application checks. No repeated cross-validation, latency benchmark, aggregate Trust Score distribution or real-world accuracy is claimed.

\section{Dataset Characteristics}
Table 7.1 records observed product counts. Both CSVs have no empty cells or duplicate full rows in the inspected snapshot; the class distribution is plotted in Figure 5.1.

\begin{longtable}{p{0.184\linewidth}p{0.184\linewidth}p{0.184\linewidth}p{0.184\linewidth}p{0.184\linewidth}}
\caption{Generated dataset summary}\\

Product & Records & Total columns & Approved & Rejected\\
Home loan & 3,000 & 18 & 1,529 & 1,471\\
Personal loan & 3,000 & 17 & 1,563 & 1,437\\
\end{longtable}
The generator defines the dependencies used for prediction, including a deliberate borderline gender penalty. Near-balanced classes and high discrimination should therefore be interpreted as outcomes within the designed simulation, not estimates for genuine bank customers.

\section{Model Performance}
Table 7.2 reports saved product-specific held-out metrics. Both test subsets contain 600 records, drawn separately from their respective datasets using the same split procedure.

\begin{longtable}{p{0.153\linewidth}p{0.153\linewidth}p{0.153\linewidth}p{0.153\linewidth}p{0.153\linewidth}p{0.153\linewidth}}
\caption{Saved product held-out metrics}\\

Product & Accuracy & Precision & Recall & F1 & ROC-AUC\\
Home loan & 0.9700 & 0.9615 & 0.9804 & 0.9709 & 0.9978\\
Personal loan & 0.9733 & 0.9744 & 0.9744 & 0.9744 & 0.9963\\
\end{longtable}
Home accuracy/ROC-AUC were reproduced as 0.97/0.9978102352051932; personal values as 0.9733333333333334/0.9963375672095379. These match stored files. The other table values are saved metrics rather than a new Phase 2 experiment. Root legacy metadata reports a different accuracy near 99.17\%; it is not used for the two current product models.

The high values are consistent with explicit learnable label rules. External data shift, observed repayment outcomes, reject inference, probability calibration and temporal validation are not evaluated by these synthetic splits.

\section{Interpretation of Classification and Fairness Metrics}
Accuracy is the fraction of correct synthetic labels. Precision concerns correctness among predicted approvals; recall concerns detection of actual approved labels; F1 combines precision and recall. ROC-AUC measures ranking across thresholds. None alone measures financial harm or calibrated probability.

Saved female/male predicted approval-rate ratios are 0.905840 for home and 0.864819 for personal loans. These describe group outcomes on the generated test subsets. They are distinct from per-request R, which compares one applicant's probabilities after gender flip. The report does not infer compliance or population fairness from either statistic [10].

\section{Worked Home-Loan Evaluation}
A retained Phase 1 HTTP case provides a traceable worked example rather than an aggregate experiment. The sample is synthetic: female, age 35, married, one dependent, graduate, salaried, ten years of experience; applicant income 80,000, co-applicant income 20,000, existing EMI 6,000; requested amount 3,500,000 over 240 months; property value 5,000,000 in an urban area; CIBIL 810 and zero missed payments.

The response was Approved with approval/predicted-class probability 0.9994001389. At 8.5\%, proposed EMI was 30,373.81 and total debt 36,373.81. FOIR was 0.3637 (36.4\% displayed), classified MODERATE. Debt-payment surplus was 63,626.19. The same profile's approval does not itself authorize a loan or verify applicant documents.

Table 7.3 records its HTTP diagnostic values. These are not mixed with direct-dictionary evaluations, whose different default expansion can alter the seed.

\begin{longtable}{p{0.46\linewidth}p{0.46\linewidth}}
\caption{Retained worked-case diagnostic values}\\

Metric & Retained HTTP value\\
Faithfulness F & 0.8448\\
Stability S & 0.9714285714\\
Consistency C & 0.5000\\
Gender-flip score R & 0.9997223616\\
Composite T & 0.7885093790\\
\end{longtable}
The arithmetic is 0.35(0.8448)+0.25(0.9714285714)+0.30(0.5)+0.10(0.9997223616)=0.7885093790. The moderate trust band coexists with very high predicted-class probability: the two fields describe different properties.

\section{Paired SHAP and LIME Example}
Figures 7.1 and 7.2 plot seven leading absolute contributions from that same retained HTTP response. Charts are newly rendered from recorded values, not screenshots or newly measured aggregate experiments. Numeric source values and provenance are retained in evidence/phase1\_retained.json.

\begin{figure}[p]\centering\includegraphics[width=\linewidth,height=.82\textheight,keepaspectratio]{\detokenize{assets/shap_case.png}}\caption{Leading TreeSHAP contributions for the retained home-loan HTTP case; raw model-margin scale.}\end{figure}
CIBIL is the strongest positive SHAP contributor (4.510115), followed by FOIR (0.956479), missed payments (0.951525) and LTV (0.838448). Gender\_Male contributes -0.573078 for this female profile. A near-one gender-flip probability score can coexist with a visible raw-margin attribution because a large margin can saturate approval probability; the quantities answer different questions.

\begin{figure}[p]\centering\includegraphics[width=\linewidth,height=.82\textheight,keepaspectratio]{\detokenize{assets/lime_case.png}}\caption{Leading approval-class LIME weights for the same retained HTTP case; local-surrogate scale.}\end{figure}
LIME's leading positive weights include CIBIL (0.578389), missed payments (0.213624) and LTV (0.095734). Both top-three sets contain CIBIL and missed payments; their union has four names, giving Jaccard consistency 2/4=0.5. The magnitudes are not compared directly to SHAP log-odds. The API does not return LIME fit quality, so agreement is not a standalone surrogate-validity test.

\section{Financial Solvency and Trust Results}
Two executed HTTP cases demonstrate the separation between statistical classification and affordability assessment.

\begin{longtable}{p{0.115\linewidth}p{0.115\linewidth}p{0.115\linewidth}p{0.115\linewidth}p{0.115\linewidth}p{0.115\linewidth}p{0.115\linewidth}p{0.115\linewidth}}
\caption{Executed approved/insolvent cases}\\

Product case & ML class & Proposed EMI & Existing EMI & Income & FOIR & Solvency & Trust\\
Home conflict & Approved & 22,303.06 & 14,000 & 67,900 & 0.5347 & INSOLVENT & 0.800388\\
Personal conflict & Approved & 4,954.55 & 13,700 & 32,600 & 0.5722 & INSOLVENT & 0.965256\\
\end{longtable}
A high Trust Score can describe the implemented explanation checks for a classification that exceeds the project's affordability guide. The personal case is particularly instructive: high trust does not constitute solvency or lending approval. The frontend recommendation can flag decline while preserving the API class.

A scan of all dataset rows found 308 home and 53 personal model approvals above 50\% FOIR. These counts cover the full generated datasets, not held-out error rates or independently sampled populations. They illustrate that the learned label rule and the online 50\% guide differ.

\section{Application and Report Presentation}
The implemented navigation supports product forms, officer summary/audit, applicant assessment, history and batch modals, submitted inputs and PDF actions. Component tests and HTTP serving provide partial interface evidence, but no interactive browser walkthrough was completed. Accordingly, no browser screenshots are fabricated or included in this draft.

Figure 7.3 is an application-generated PDF specimen re-rendered from the retained worked-case input and result. It demonstrates report presentation, not fresh model inference. The reference identifies the retained evaluation; the generation date belongs to this document-rendering operation.

\begin{figure}[p]\centering\includegraphics[width=\linewidth,height=.82\textheight,keepaspectratio]{\detokenize{assets/application_report.png}}\caption{Application-generated PDF specimen using the retained home-loan evaluation; not a browser screenshot.}\end{figure}
The specimen shows the financial summary and weighted pillars. Its template contains provisional sanction language and signature placeholders. As discussed in Section 5.20, static narrative claims must not be interpreted as independently verified financial facts. Browser captures can be added only when genuine assets are available.

\section{Discussion of Findings}
The prototype connects product prediction, dual explanation, diagnostics, affordability assessment, storage and reporting in a usable engineering workflow. Saved classification metrics and recomputed discrimination values agree for the synthetic experiment. Worked cases demonstrate why class probability, explanation diagnostics and repayment burden should remain separate.

The tests also reveal operational boundaries: incomplete financial-domain validation, stale frontend/batch expectations, mixed display-value semantics and an optimistic fairness error fallback. These observations do not erase the implemented workflow; they define what the evidence supports. There is no controlled aggregate explanation benchmark or real-lender deployment claim. The contribution is integration and transparent inspection of the implemented diagnostics.

\section{Summary}
Product models reached approximately 97.00\% and 97.33\% accuracy on their generated test records. The paired example demonstrates distinct SHAP/LIME scales, top-feature agreement and weighted trust arithmetic. Approved/insolvent cases establish advisory solvency behavior. This combination supports a prototype demonstration and a reproducible discussion, while broader lending validity remains future evaluation.

\chapter{Project Planning}
\section{Development Methodology}
Repository history supports an iterative chronology rather than a formal sprint baseline. Table 8.1 retains the original date windows and associates them with recorded changes. Commit dates identify recorded integration activity, not continuous effort or exact completion times.

\begin{longtable}{p{0.307\linewidth}p{0.307\linewidth}p{0.307\linewidth}}
\caption{Git-supported development chronology}\\

Recorded window & Evidence / representative commit & Activity\\
8-11 September 2026 & 839480b; 2a5d282; e04a93b & Initial project, deterministic explanation evaluation and UI/roadmap work.\\
29 September 2026 & bfadff4; 99c796e & Assessment wording and audit explanation revision.\\
17 October 2026 & f280c08; 64ba516 & UI simplification, tests, sizing and presentation material.\\
22 October 2026 & 416deb2; 6a27e18; aea9b02; 4807ae4 & History, batch CSV, schema initialization and test additions.\\
24 October 2026 & d807962 & Documentation revisions.\\
30 October-1 November 2026 & e231a48; d2d0c79; 82fda67; 6366570; 5a7c1ff & Product refinements, trust scoring, SQLite, PDF/calculator and navigation changes.\\
\end{longtable}
The available history begins on 8 September and ends on 1 November 2026. This chapter does not infer individual contributions from repository authorship.

\section{Project Timeline and Milestones}
The available Git history begins on 8 September 2026 and includes project changes through 1 November 2026. The named milestones in Table 8.1 are based on commit dates and subjects. They do not measure effort or prove that work was finished on the commit date. A formal week-by-week Gantt schedule was not located in the repository.

\section{Team Membership and Contribution Records}
The report identifies Ashwin Patil (46), Aaditya Patil (44), Kanad Nerkar (41) and Ved Kalantri (28) as the project group. These names and roll numbers are retained from the supplied report. The final individual contribution allocation requires a group-confirmed record; no engineering responsibility is invented from commit authors or filenames. The present draft records collective project activities and leaves individual allocation to that separate confirmation.

\section{Challenges}
Implementation challenges reflected in the code include keeping two product feature sets and preprocessing pipelines aligned, mapping explainers to transformed model features, making stochastic LIME output reproducible for repeated input, maintaining local database operation when PostgreSQL is not available, and separating prediction results from explanationquality indicators. The code and tests also expose a point requiring review: some test wording describes an automatic solvency rejection while the route currently returns model prediction and solvency assessment as separate results.

Verification later exposed serialization-version warnings, invalid-input acceptance and stale regression fixtures. These are observed technical findings. Whether they were experienced as development challenges by particular members requires the group's own account.

\section{Chronological Activity Chart}
Table 8.2 is a milestone chronology, not a duration-based Gantt chart. The repository does not provide start/end estimates or an approved task-duration schedule. A formal Gantt can be included only if the group supplies its genuine plan.

\begin{longtable}{p{0.131\linewidth}p{0.131\linewidth}p{0.131\linewidth}p{0.131\linewidth}p{0.131\linewidth}p{0.131\linewidth}p{0.131\linewidth}}
\caption{Chronological activity chart}\\

Activity & Sep 8-11 & Sep 29 & Oct 17 & Oct 22 & Oct 24 & Oct 30-Nov 1\\
Initial architecture, XAI and UI & Recorded &  &  &  &  & \\
Assessment/audit wording &  & Recorded &  &  &  & \\
Testing and UI refinements &  &  & Recorded &  &  & \\
Database, history and batch &  &  &  & Recorded &  & \\
Documentation &  &  &  &  & Recorded & \\
Product/final refinements &  &  &  &  &  & Recorded\\
\end{longtable}
The dated evidence makes project evolution visible without treating isolated commits as measured development durations.

\chapter{Conclusion and Future Scope}
\section{Conclusion}
TrustFin AI implements a loan decision-support workflow that links home-loan and personalloan models with SHAP and LIME explanations, four explanation checks, a weighted Trust Score, a React dashboard, database persistence, batch evaluation, and PDF reporting. The product models use generated synthetic records and show high stored metrics on stratified hold-out splits. These results are useful for demonstrating the system workflow, but they do not establish real-world underwriting performance.

The project makes an engineering contribution by presenting the prediction and explanation checks as separate outputs within one application. It also makes several assumptions visible: the faithfulness score is heuristically calibrated, stability depends on a specified perturbation, consistency compares a top-three feature set, and fairness is measured through an individual gender-flip probability difference. The confidence value is the predicted-class probability and is not independently calibrated. These distinctions support an accurate interpretation of the prototype system.

The local checks verified selected prediction, explanation, history, batch and reporting operations while retaining failed tests and untested browser paths. The main practical value is a reviewable assessment record that separates classifier output, explanation diagnostics and a financial affordability guide.

\section{Project Contributions}
The implemented contributions are:

• two product-specific feature and preprocessing pipelines for home and personal loans; • XGBoost training and saved test metrics for generated datasets; • local TreeSHAP and tabular LIME explanation paths; • code-level metrics for faithfulness, stability, top-feature consistency, and individual counterfactual sensitivity; • a weighted Trust Score and a separate predicted-class probability field; • a web interface, evaluation and batch history, local SQLite fallback, and on-demand PDF reporting.

\section{Limitations}
The principal limitations are synthetic-only data, limited external validation, no probability calibration, no group-level fairness audit in the request path, no causal explanation guarantees, and a custom composite score without external validation. The application does not implement authentication or server-enforced role access. The what-if calculator changes affordability inputs but does not request a fresh model prediction. The project also has no stored product-specific confusion matrices or held-out aggregate Trust Score study.

Additional observed boundaries are advisory-only solvency, incomplete financial-domain validation, a fairness exception fallback of 1.0, transformed-value display inconsistencies and static report narratives. These define the current prototype evidence rather than hypothetical defects in a production installation.

\section{Future Scope}
Future work should first evaluate the models on a lawful, representative, de-identified lending dataset and establish leakage controls, missing-data policy, calibration, subgroup metrics, and temporal validation. Explanation metrics should be tested against independently justified perturbation schemes and human review, with sensitivity analyses for thresholds and weights. The Trust Score should be treated as a hypothesis until an evaluation study demonstrates how it relates to meaningful explanation quality or reviewer decisions.

The application could add authentication, server-side authorization, privacy controls, audit retention policies, versioned database migrations, secure deployment configuration, monitoring, and clear human-review workflows. The what-if calculator could call a controlled scenario endpoint if the intended behavior is to recompute model output; such a change would require displaying the changed inputs and preventing the scenario from being mistaken for an approved loan. Additional explainers or constrained counterfactual methods may be evaluated as future research, not described as existing functions.

Live credit-bureau integration is a relevant extension subject to authorized access and validated data mapping. The existing financial what-if calculator can evolve into model-verified interactive simulation with feasible-change constraints. Portfolio-level fairness auditing can add subgroup metrics, uncertainty and monitoring beyond individual gender sensitivity. These extensions belong to future work and are not counted as completed prototype outcomes.

\section{Closing Statement}
The prototype implementation brings loan prediction, local explanation, explanation validation, history, batch processing, and reporting into one decision-support system. Its contribution is the integration of those functions with explicitly defined diagnostics and an interface for reviewing results. Extension to real lending operations would require institution-specific data governance, security controls, external model validation, and human evaluation.

\chapter*{Bibliography}\addcontentsline{toc}{chapter}{Bibliography}
[1] D. J. Hand and W. E. Henley, Statistical classification methods in consumer credit scoring: A review, Journal of the Royal Statistical Society: Series A, 160(3), 523-541, 1997. DOI: 10.1111/j.1467-985X.1997.00078.x.

[2] S. Lessmann, B. Baesens, H.-V. Seow and L. C. Thomas, Benchmarking state-of-the-art classification algorithms for credit scoring: An update of research, European Journal of Operational Research, 247(1), 124-136, 2015. DOI: 10.1016/j.ejor.2015.05.030.

[3] C.-K. Yeh, C.-Y. Hsieh, A. Suggala, D. I. Inouye and P. K. Ravikumar, On the (In)fidelity and Sensitivity of Explanations, Advances in Neural Information Processing Systems 32, 2019. Official record: https://proceedings.neurips.cc/paper/2019/hash/a7471fdc77b3435276507cc8f2dc2569-Abstract.html.

[4] D. Alvarez-Melis and T. S. Jaakkola, On the Robustness of Interpretability Methods, arXiv preprint 1806.08049, 2018. DOI: 10.48550/arXiv.1806.08049.

[5] E. Tabassi, Artificial Intelligence Risk Management Framework (AI RMF 1.0), NIST AI 100-1, National Institute of Standards and Technology, 2023. DOI: 10.6028/NIST.AI.100-1.

[6] J. H. Friedman, Greedy function approximation: A gradient boosting machine, The Annals of Statistics, 29(5), 1189-1232, 2001. DOI: 10.1214/aos/1013203451.

[7] T. Chen and C. Guestrin, XGBoost: A Scalable Tree Boosting System, Proceedings of the 22nd ACM SIGKDD International Conference on Knowledge Discovery and Data Mining, 785-794, 2016. DOI: 10.1145/2939672.2939785.

[8] S. M. Lundberg and S.-I. Lee, A Unified Approach to Interpreting Model Predictions, Advances in Neural Information Processing Systems 30, 4765-4774, 2017. Official record: https://proceedings.neurips.cc/paper/2017/hash/8a20a8621978632d76c43dfd28b67767-Abstract.html.

[9] M. T. Ribeiro, S. Singh and C. Guestrin, Why Should I Trust You?: Explaining the Predictions of Any Classifier, Proceedings of the 22nd ACM SIGKDD International Conference on Knowledge Discovery and Data Mining, 1135-1144, 2016. DOI: 10.1145/2939672.2939778.

[10] M. Hardt, E. Price and N. Srebro, Equality of Opportunity in Supervised Learning, Advances in Neural Information Processing Systems 29, 2016. Official record: https://proceedings.neurips.cc/paper\_files/paper/2016/hash/6a9659feb1216f14f7384ba499518b38-Abstract.html.

[11] S. Wachter, B. Mittelstadt and C. Russell, Counterfactual Explanations without Opening the Black Box: Automated Decisions and the GDPR, Harvard Journal of Law \& Technology, 31(2), 841-887, 2018. Publisher PDF: https://jolt.law.harvard.edu/assets/articlePDFs/v31/Counterfactual-Explanations-without-Opening-the-Black-Box-Sandra-Wachter-et-al.pdf.

[12] S. M. Lundberg et al., From local explanations to global understanding with explainable AI for trees, Nature Machine Intelligence, 2, 56-67, 2020. DOI: 10.1038/s42256-019-0138-9.

\chapter*{Appendix A - Observed Verification Environment}\addcontentsline{toc}{chapter}{Appendix A - Observed Verification Environment}
The retained Phase 1 environment is recorded for reproducibility. Versions are observations rather than a dependency upgrade recommendation.

\begin{longtable}{p{0.46\linewidth}p{0.46\linewidth}}
\caption{Observed verification environment}\\

Component & Observed version\\
Python & 3.12.7\\
Node.js / npm & 22.19.0 / 10.9.3\\
FastAPI / Uvicorn & 0.162.0 / 0.54.0\\
Pydantic / SQLAlchemy & 2.14.2 / 2.1.6\\
pandas / NumPy & 3.1.0 / 2.5.4\\
scikit-learn / XGBoost & 1.9.1 / 3.4.1\\
SHAP / LIME & 0.54.1 / 0.2.0.1\\
ReportLab / pytest & 5.1.0 / 9.2.0\\
Vitest & 5.0.5\\
\end{longtable}
Backend verification used in-memory SQLite and existing model artifacts. Browser automation failed before opening the page. No real applicant records were introduced. The host hardware specification was not captured in the retained evidence.

\chapter*{Appendix B - Implementation and Evidence Map}\addcontentsline{toc}{chapter}{Appendix B - Implementation and Evidence Map}
\begin{longtable}{p{0.46\linewidth}p{0.46\linewidth}}
\caption{Source and evidence map}\\

Responsibility & Principal source or evidence\\
API entry and validation errors & backend/app/main.py\\
Prediction/batch/history/report routes & backend/app/api/routes/prediction.py\\
Request/response schemas & backend/app/api/schemas.py\\
Product loading and defaults & backend/app/ml/registry.py\\
Financial transformations & backend/app/ml/preprocessing\_homeloan.py and preprocessing\_personalloan.py\\
Product training & scripts/train\_homeloan.py and train\_personalloan.py\\
Synthetic generation & scripts/generate\_homeloan\_dataset.py and generate\_personalloan\_dataset.py\\
Local explainers & backend/app/explainability/shap\_explainer.py and lime\_explainer.py\\
Pillars and composite & backend/app/validation/*.py; backend/app/trust/trust\_score.py\\
Persistence & backend/app/db/models.py and session.py\\
Application report & backend/app/reporting/pdf\_generator.py\\
React navigation & frontend/src/App.jsx\\
Recommendations/scenarios & frontend/src/utils/decisionAdvisory.js; components/RecourseSimulator.jsx\\
Software tests & backend/tests; frontend/src/test\\
Stored metrics & models/homeloan/metrics.json; models/personalloan/metrics.json\\
Draft evidence snapshots & evidence/model\_metrics\_snapshot.json; evidence/phase1\_retained.json\\
\end{longtable}
Selected backend module results, extra batch checks and frontend totals in Chapter 6 are retained from the Phase 1 execution record. They are not represented as newly executed full-suite logs. Source-integrity hashes are recorded alongside the draft assets.

\chapter*{Appendix C - Worked-Case Reproducibility Record}\addcontentsline{toc}{chapter}{Appendix C - Worked-Case Reproducibility Record}
The worked case in Sections 7.5-7.6 uses the same retained HTTP evaluation for both explanation charts. The JSON record includes the synthetic input, probability, financial values, four scores and selected contribution values. It is a curated evidence excerpt, not a complete aggregate experiment or an untouched raw terminal log.

Approval probability was 0.9994001388549805. Faithfulness was 0.8448, stability 0.9714285714285714, consistency 0.5 and gender sensitivity 0.9997223615646362. Weighted trust was 0.7885093790136065. The report specimen is generated from these retained values without re-evaluating the model. Original request defaults and computational environment must match to reproduce applicant-derived random seeds.

\end{document}